%% CPP 2027 Submission — 0pat Exercise XXVII: Split-Complex Numbers (Imaginary-Radius Orthogonal Euler Circles)
\documentclass[sigplan,10pt,anonymous,review]{acmart}
\settopmatter{printfolios=true,printacmref=false}

\AtBeginDocument{%
  \providecommand\BibTeX{{\normalfont B\scshape ib\TeX}}}

\settopmatter{printacmref=false}
\renewcommand\footnotetextcopyrightpermission[1]{}
\pagestyle{plain}

\usepackage{microtype}
\usepackage{booktabs}
\usepackage{amsmath}
% XeTeX (tectonic): acmart loads newtxmath/libertine which already defines \Bbbk;
% reset it so amssymb's definition is accepted.
\let\Bbbk\relax
\usepackage{amssymb}

\begin{document}

\title{Split-Complex Numbers as Imaginary-Radius Orthogonal Euler
Circles (Exercise XXVII)}

\author{Anonymous}
\affiliation{%
  \institution{Anonymous}
}
\email{anonymous@example.org}

\begin{abstract}
We report the formal implementation of an observation frame suggested by
the pat exercise series: two orthogonal Euler circles, one of radius $1$
and one of radius $i$ --- an imaginary-radius circle, which cannot exist
in the complex plane (radii are nonnegative reals, $|z| = i$ has no
solution). The host structure is the split-complex numbers
$a + bj$, $j^2 = +1$, with pseudo-norm $\|a + bj\|^2 = a^2 - b^2$
(Minkowski metric): the radius-$1$ circle is the hyperbola
$a^2 - b^2 = 1$, and the radius-$i$ circle is the conjugate hyperbola
$a^2 - b^2 = -1$, nonempty (its unit $j$ lies on it). Seven theorems are
machine-checked in Lean 4 / mathlib (with a self-built 30-line
split-complex structure, standard mathematics): $j^2 = 1$,
pseudo-norm multiplicativity $\|zw\|^2 = \|z\|^2\|w\|^2$, nonemptiness
of both circles, the orthogonality condition $1^2 + i^2 = 0$, disjointness
of the conjugate hyperbolas, and the observation contrast: every odd
prime is a split pseudo-norm $p = ((p+1)/2)^2 - ((p-1)/2)^2$, whereas in
the complex plane only primes $p \equiv 1, 2 \pmod 4$ are norms (Fermat).
Zero errors, zero \texttt{sorry}.
\end{abstract}

\keywords{Lean, formalization, split-complex numbers, hyperbola, imaginary radius, 0pat}

\maketitle

\section{The observation frame}

The exercise series has been observing the prime axis through circles:
the circle $\mathbb{Z}/p\mathbb{Z}$ (Exercise XXIV), the orthogonal
coordinate frame of two circles (Exercise XXV), and the critical circle
of the Riemann direction. The question posed here is whether the two
orthogonal Euler circles may carry \emph{different units} --- one of
radius $1$ (natural unit) and one of radius $i$ (imaginary unit) ---
without forcing the Cartesian constraint of the complex plane, where the
real and imaginary axes share the same unit.

The answer is that such a frame cannot live in the complex plane: a
circle $\{z : |z - c| = r\}$ requires $r \ge 0$, and $|z| = i$ has no
solution. The host structure is the \emph{split-complex} numbers, the
two-dimensional real algebra $a + bj$ with $j^2 = +1$ (contrast
$\mathbb{C}$, where $J^2 = -1$). Its pseudo-norm is the Minkowski
metric $\|a + bj\|^2 = a^2 - b^2$. In this structure, the radius-$1$
circle is the hyperbola $a^2 - b^2 = 1$, and the radius-$i$ circle is
the conjugate hyperbola $a^2 - b^2 = -1$, which is nonempty --- its unit
$j$ lies on it. The orthogonality of the two radii is exact in the
split metric: $r_1^2 + r_2^2 = 1^2 + i^2 = 1 + (-1) = 0$, the condition
for orthogonal concentric circles.

\section{The theorems}

Seven theorems, all machine-checked with zero \texttt{sorry}. The
split-complex structure is defined in 30 lines (a pair of reals with
multiplication $(a,b)(c,d) = (ac + bd, ad + bc)$); the theorems are
standard mathematics.

\paragraph{The defining algebra (j\_sq).} $j^2 = 1$. This is the
single structural difference from the complex numbers ($J^2 = -1$) and
the origin of every subsequent property.

\paragraph{Pseudo-norm multiplicativity (normSq\_mul).}
$\|zw\|^2 = \|z\|^2\|w\|^2$. The split pseudo-norm behaves like the
complex norm under multiplication; the proof is a direct expansion.

\paragraph{Nonemptiness of both circles.} $\|j\|^2 = -1$ (the
radius-$i$ circle contains its unit $j$) and $\|1\|^2 = 1$ (the
radius-$1$ circle contains $1$). The imaginary-radius circle exists and
is inhabited --- the property that fails in $\mathbb{C}$.

\paragraph{Orthogonality (orthogonal\_radii).} $1 + (-1) = 0$: in the
split metric, the radii $1$ and $i$ satisfy $r_1^2 + r_2^2 = 0$, the
concentric-circle orthogonality condition. The two Euler circles are
orthogonal in exactly the sense the frame requires.

\paragraph{Disjointness (circles\_disjoint).} No point lies on both
circles: $a^2 - b^2 = 1$ and $a^2 - b^2 = -1$ cannot hold together. The
conjugate hyperbolas are disjoint --- the frame has two separate loci,
which is what makes them a coordinate pair rather than one circle.

\paragraph{Observation contrast (odd\_prime\_is\_split\_norm).}
Every odd prime $p$ is a split pseudo-norm:
$p = ((p+1)/2)^2 - ((p-1)/2)^2$. In the complex plane, only primes
$p \equiv 1, 2 \pmod 4$ are norms (Fermat's two-square theorem, Exercise
XXIV). The split frame therefore observes primes without discrimination
--- every odd prime is representable --- whereas the complex frame
distinguishes half of them. The contrast is the frame's first
observational result: the choice of unit ($1$ vs $i$) determines the
discriminating power of the observation.

\section{Honest boundary and position}

As in every exercise: the content is classical mathematics (the
split-complex numbers are a standard structure; all theorems are
known), the proofs are machine-checked, and no claim is made beyond the
statements. The value is positional: the frame requested by the series
(radius-$1$ and radius-$i$ orthogonal Euler circles) is now implemented
and its basic observational property (no prime discrimination) is
established. The frame is available for further observation --- for
example, of structures that the complex plane cannot host.

\section{Formalization}

The proof is a single Lean file (\texttt{proof.lean}, 135 lines),
importing only \texttt{Mathlib.Data.Real.Basic},
\texttt{Mathlib.Data.Nat.Prime.Basic}, \texttt{Mathlib.Tactic.Ring}, and
\texttt{Mathlib.Tactic.Linarith}. The split-complex structure
(\texttt{SplitComplex}, 30 lines) is a self-contained standard
definition; no pat framework concepts are used. Compiled with Lean
4.32.2 / mathlib v4.32.2, zero errors, zero \texttt{sorry}.

\begin{center}
\begin{tabular}{lll}
\toprule
Theorem & Statement & Proof core \\
\midrule
\texttt{j\_sq} & $j^2 = 1$ & ext; simp \\
\texttt{normSq\_mul} & $\|zw\|^2 = \|z\|^2\|w\|^2$ & simp; ring \\
\texttt{imaginary\_radius\_circle} & $\|j\|^2 = -1$ & simp \\
\texttt{real\_radius\_circle} & $\|1\|^2 = 1$ & simp \\
\texttt{orthogonal\_radii} & $1 + (-1) = 0$ & norm\_num \\
\texttt{circles\_disjoint} & $\neg\exists z,\ \|z\|^2 = 1 \wedge \|z\|^2 = -1$ & linarith \\
\texttt{odd\_prime\_is\_split\_norm} & $p = ((p{+}1)/2)^2 - ((p{-}1)/2)^2$ & div\_add\_mod; omega \\
\bottomrule
\end{tabular}
\end{center}

\section{Conclusion}

The imaginary-radius orthogonal Euler frame is implemented: split-complex
numbers host the radius-$i$ circle, the orthogonality condition
$1^2 + i^2 = 0$ holds exactly, and the frame's first observation is that
it does not discriminate primes (every odd prime is a pseudo-norm),
in contrast to the complex frame. The frame now exists, machine-checked,
ready for further observation.

\bibliographystyle{ACM-Reference-Format}
\bibliography{refs}

\end{document}
